\chapter{RESEARCH METHODOLOGY}
\label{ChapterResearchMethodology}

\section{Research Approach}
We use a hybrid research design combining the principles of design
science research and applied research methodology in implementing
this study. The main purpose of the work is to design, develop, and
evaluate an intelligent program which analyzes Terms of Service (ToS)
documents and helps detect potentially risky clauses. The research of
design science is evidenced by an operative artifact, which can be
known as RiskWise, that uses the latest techniques of Natural
Language Processing (NLP), Knowledge Graphs, and Retrieval-Augmented
Generation (RAG) for solving a real-world problem. Simultaneously, we
see the elements of applied research, as the project aspires for a
practical technological solution to the challenge of user
comprehension of complex legal agreements. Research involved
iterative problem understanding, literature exploration, system
design, implementation, and validation through both experimental and
case-based testing. This allowed for iterative improvements in system
architecture and analysis pipeline based on insights obtained during
development.

\section{Preliminary Research and Literature Exploration}
We did a preliminary literature review to get an idea of the current
research on things like AI-assisted contract understanding, unfair
clause detection, and legal document analysis. We acquired pertinent
research articles and technical materials from academic databases and
repositories (e.g., Google Scholar, arXiv, IEEE Xplore) via reference
tracing from specific publications. The literature search focused on
keywords such as Terms of Service clause detection, unfair ToS
detection utilising NLP, legal document analysis, knowledge
graph-based legal reasoning, Retrieval-Augmented Generation for legal
documents, and Graph-RAG architectures. We looked at about 10 to 15
papers and chose 6 to 8 to read closely because they were new,
relevant to legal document processing, methodologically innovative,
and useful for explainable AI-driven legal analysis.
\\
This exploratory phase helped us find important research areas that
need more study, such as the fact that traditional NLP models are not
good at structural reasoning, there are no clear ways to identify
risks, and there is a need for hybrid search strategies that combine
semantic similarity with relational reasoning.

\section{System Development Methodology}
The RiskWise system was made using a method that involved making
prototypes and trying things out over and over. The project goes
through a number of iterations. It starts with exploratory exercises
using different clause detection models and ends with the design of a
hybrid architecture that includes building a knowledge graph and
RAG-based reasoning.
\\
First, we did comparative analyses of transformer-based models for
clause extraction tasks. This helped us understand the pros and cons
of text-based methods. Based on these results, the system
architecture was gradually improved to include structured ways to
store legal information and ways to find it that take into account the context.
The iterative method made the system's parts modular, such as the
document ingestion, semantic chunking, embedding generation,
knowledge graph construction, hybrid retrieval logic, and generative
reasoning modules. Every release of the work gradually improved the
functionality of the features without changing how they worked with
earlier versions. This way, the system's abilities can gradually improve.

\section{Proposed Methodological Framework}
The methodological framework of the project is based on a multi-stage
processing pipeline designed to transform unstructured legal
documents into an interpretable knowledge representation and risk
analysis output. The major methodological stages are summarized below.

\section{Document Ingestion and Semantic Parsing}
It starts when users upload ToS documents in supported formats. Text
extraction methods are used to extract raw textual content from
documents, and semantic segmentation is performed into meaningful
chunks. Semantic chunking\index{Semantic Chunking}, as opposed to fixed-length segmentation
that discards contextual boundaries such as clauses and sections, can
enhance the quality of downstream retrieval.

\subsection{Embedding Generation and Vector Indexing}
Semantic embeddings\index{Semantic Embedding} are obtained for each segment of the text chunk,
based on our domain-adapted embedding models which can represent
legal semantics. Similarity-based retrieval is performed by storing
such embeddings in a vector index, which facilitates efficient
similarity-based retrieval during query processing. The vector\index{Vector Database}
component of the system is responsible for fuzzy matching and
semantic equivalence detection of legal terminology variations.

\subsection{Knowledge Graph Construction}
Alongside embedding generation, the system builds out a knowledge
graph for the document. Large Language Models\index{LLM} are used to derive
subject-predicate-object triples\index{Triple Extraction} that describe relationships between
entities such as service providers, users, data, and obligations.
These triples are mapped to a structured legal ontology and stored in
a graph database\index{Graph Database}, preserving relational and logical dependencies
within the document.

\subsection{Hybrid Retrieval Mechanism}
In the system we implement a hybrid retrieval strategy\index{Hybrid Retrieval} involving
vector similarity\index{Vector Similarity Search} and knowledge graph traversal\index{Graph Traversal}. Vector retrieval
retrieves semantically relevant clauses, and graph traversal finds
relationally connected information which don't share the common
vocabulary. Incorporating these retrieval signals forms the
comprehensive background information that further justifies the
reasoning ability.

\subsection{Generative Reasoning and Risk Analysis}
A generative language model\index{Generative Reasoning} set up as a legal risk assessor\index{Risk Analysis} gets the
combined contextual information. The model analyses clauses, finds
potentially risky parts of contracts, gives them severity scores, and
makes natural language summaries that are easy to understand based on
the context it finds. Using citations to prompt further increases the
transparency and traceability of outputs.

\section{Evaluation Strategy}
The proposed system underwent a system evaluation utilising
qualitative case-study analysis. We looked at a few Terms of Service
documents with the developed pipeline to see how it works, how well
it can find risks, and how well it can respond. Experimental phases
culminating in comparative model evaluation have yielded additional
empirical evidence for baseline clause extraction performance. A
manual inspection of system outputs was conducted to confirm the
relevance of identified clauses, the appropriateness of risk
categorization, and contextual accuracy of generated explanations.
Since the project was exploratory and prototype-based, it was thought
that a qualitative validation method would be best.

\section{Validation and Reliability Measures}
To ensure the system's results are reliable and trustworthy, a method
has been set up that includes several validation steps.
Retrieval-Augmented Generation allowed the model to build its answers
based on a document it had found instead of just using the knowledge
it had already learned. Knowledge graph grounding made sure that
clauses were consistent with each other, which made reasoning even
harder. Other ways to check reliability included structured parsing
of the information that was taken out, fallback mechanisms for
dealing with model outputs that weren't consistent, and using
supporting text segments in context. This combined method lowers the
chance of hallucination and also makes it easier to understand the
results of system-generated analysis.
